\documentclass[11pt]{article}

% Compile from the repository root with:
%   tectonic docs/scaling_imperfect_reward_report.tex
% Figures come from:
%   uv run python experiments/run_scaling_states.py --seeds 3
%   uv run python experiments/run_scaling_actors.py --seeds 5
%   uv run python experiments/run_scaling_states_actors.py --seeds 3

\usepackage[margin=1in]{geometry}
\usepackage{amsmath, amssymb}
\usepackage{booktabs}
\usepackage{float}
\usepackage{graphicx}
\usepackage{hyperref}

\graphicspath{{./}{docs/figures/}{figures/}}
\hypersetup{colorlinks=true, linkcolor=blue, urlcolor=blue}

\setlength{\parindent}{0pt}
\setlength{\parskip}{0.6em}

\newcommand{\reportfigure}[3]{%
  \begin{figure}[H]
    \centering
    \includegraphics[width=0.86\linewidth]{#1}
    \caption{#2}
    \label{#3}
  \end{figure}
}

\title{\textbf{Scaling notes: imperfect reward (Sections 4 and 6)}}
\author{Bo Xu}
\date{}

\begin{document}
\maketitle

These are my notes on scaling the imperfect-reward side, where one Observer has a
fixed attention budget and only reports a reward back to an Actor with probability
$\theta(\mu_j) = 1 - e^{-c\mu_j}$ (it spends the budget $\mu$ on the Actors it
rates highest). I followed the ladder we agreed on: scale states first, then
actors, then both. Every comparison keeps the step count fixed, so it's at equal
compute. Short version: states alone and actors alone are fine; doing both at once
works up to a point and then falls off a cliff, and I think I know why.

\section{States alone (1 Actor, 1 Observer)}

The state is a cell on a 2-D grid and the Actor sees its coordinates, so nearby
cells have similar best actions (I avoided random per-state actions on purpose, as
you suggested). One thing tripped me up first: with all-positive rewards, plain
REINFORCE just reinforces whatever it samples first and collapses. Subtracting a
running-mean baseline fixed it, and I left it as an opt-in flag so the older cases
don't change.

After that it just works. Going from 4 up to 1024 states, the clear cells are
always right; the only misses are the thin band of cells sitting on a region
boundary, where the best action is genuinely ambiguous. Accuracy doesn't keep
dropping with size (0.90, 0.90, 0.93 at 64, 256, 1024 states) because that
boundary band is a smaller share of a finer grid.

\begin{table}[H]
\centering
\begin{tabular}{rrrr}
\toprule
states & greedy acc & mean $P(\text{correct})$ & representative $P(\text{correct})$ \\
\midrule
4    & 1.000 & 1.000 & 1.000 \\
16   & 1.000 & 0.999 & 1.000 \\
64   & 0.901 & 0.899 & 1.000 \\
256  & 0.897 & 0.893 & 1.000 \\
1024 & 0.928 & 0.921 & 1.000 \\
\bottomrule
\end{tabular}
\caption{States alone: stays high all the way to 1024; the well-inside cells stay
perfect. So state count by itself isn't the bottleneck.}
\end{table}

\section{Actors alone (1 Observer, many Actors)}

Now the opposite: one trivial state, but a growing crowd of Actors, two of which
are good and the rest are noisy distractors. The question is whether the fixed
budget can still land on the two good ones as the crowd grows.

It keeps \emph{finding} them the whole way: the two good Actors are almost always
the top-ranked ones (top-$k$ precision $\approx 1.0$, even at 256). What it loses
is the ability to put the \emph{whole} budget on them. Each bad Actor only skims a
tiny slice, but with hundreds of them the leak adds up, so the share of budget on
the good Actors drops from $\approx 0.99$ at 4 Actors to $\approx 0.59$ at 256.
Ranking survives, allocation dilutes -- which is just what the equations say,
not a bug.

\reportfigure{scaling_actors_summary}{Actors alone: the Observer keeps ranking the
good Actors at the top (top-$k$ precision flat near 1.0), but a fixed budget gets
spread thinner as the crowd grows.}{fig:actors}

\section{What happens when you increase the budget $B$?}

After sharing the actors-alone results, Peter asked whether keeping $B$ fixed
was the right choice, or whether scaling $B$ with the crowd would help. I ran
two follow-up experiments.

\textbf{Case A -- proportional scaling ($B = N/2$, 2 good Actors out of $N$).}

\begin{table}[H]
\centering
\begin{tabular}{rrrrc}
\toprule
$N$ & $B$ & conc & top-2 & $B \times \text{conc}$ \\
\midrule
4   & 2   & 1.000 & 1.00 & 2.0 \\
16  & 8   & 0.583 & 1.00 & 4.7 \\
64  & 32  & 0.166 & 0.90 & 5.3 \\
256 & 128 & 0.046 & 1.00 & 5.9 \\
\bottomrule
\end{tabular}
\caption{Scale N with $B = N/2$: concentration ratio collapses, but the
absolute reports to good Actors ($B \times \text{conc}$) actually grows.}
\end{table}

\textbf{Case B -- fixed $N = 32$, increasing $B$.}

\begin{table}[H]
\centering
\begin{tabular}{rrrr}
\toprule
$B$ & conc & top-2 & $B \times \text{conc}$ \\
\midrule
2  & 0.948 & 1.00 & 1.9 \\
4  & 0.719 & 1.00 & 2.9 \\
8  & 0.507 & 1.00 & 4.1 \\
16 & 0.344 & 1.00 & 5.5 \\
32 & 0.242 & 1.00 & 7.8 \\
\bottomrule
\end{tabular}
\caption{Fix $N = 32$, vary $B$: top-2 precision stays at 1.0 throughout, but
concentration falls monotonically with $B$.}
\end{table}

The same picture emerges in both cases: the Observer always \emph{identifies}
the good Actors correctly (top-$k$ stays near 1.0), but a larger $B$ makes it
harder to concentrate the whole budget on just two of them. The softmax has a
fixed utility gap to work with; raising $B$ does not widen that gap, so the
relative concentration shrinks. At the same time, $B \times \text{conc}$
grows, meaning good Actors do receive more absolute reward feedback -- so the
effect on learning is ambiguous. My guess is that the sharpness parameter would
need to scale with $B$ to keep concentration stable, but I have not tested that.

\section{Both at once}

Here I combine the two: the grid task (many states) \emph{and} a crowd of Actors.
Two of them can see the state and learn a real policy; the rest are ``blind''
(their input is zeroed) and can only learn one fixed action. I grow both axes
together at a fixed budget. The ``system score'' below is just accuracy weighted
by how much budget each Actor got, so it's only high when the good Actors learn
\emph{and} the Observer follows them.

\begin{table}[H]
\centering
\begin{tabular}{rrrrr}
\toprule
states & actors & good-actor acc & budget on good & system score \\
\midrule
16  & 8   & 1.000 & 0.942 & 0.956 \\
36  & 16  & 0.907 & 0.832 & 0.803 \\
64  & 32  & 0.638 & 0.531 & 0.478 \\
100 & 64  & 0.302 & 0.052 & 0.258 \\
144 & 128 & 0.262 & 0.000 & 0.250 \\
\bottomrule
\end{tabular}
\caption{Both at once: fine up to about 64 states / 32 actors, then it falls off a
cliff to chance (0.25 with 4 actions).}
\end{table}

\reportfigure{scaling_states_actors_summary}{Both at once: graceful for the first
few rungs, then everything (accuracy, budget concentration, system score) drops
together past 32 actors.}{fig:both}

\section{Why it falls off the cliff (a question for you)}

This is the one place I'm stuck and would like your read. The thing I can measure
is that the budget starts spread evenly, $\mu_j = B/N$, so the chance an Actor's
reward gets reported back is $1 - e^{-cB/N}$ -- about 0.39 at 8 Actors but only
0.03 at 64. So at large $N$ a good Actor gets very few reports early on, and it
looks like its utility never gets a chance to separate from the blind crowd.

What I don't understand is why it collapses so sharply -- to chance -- rather than
just degrading more gently like the actors-alone case did. Is it the effect I
suspect (too few early reports, so nobody learns and nobody gets followed), or is
something else going on? And how would you suggest getting around it? I didn't want
to start trying fixes before checking with you what the right way to think about it
is.

\section{Where this leaves us}

States alone: fine to 1024. Actors alone: ranking stays correct to 256, budget
just dilutes. Both together: fine to about 64 states / 32 actors, then the cliff.

Everything's committed with tests, and there are longer per-phase notes in the
repo if you want the details.

\end{document}
